\documentclass{jsarticle}
\usepackage{amsmath}
\usepackage{amssymb}
\begin{document}
\title{}
\author{}
\date{}
\maketitle
  
  $$ 6, 312, {5}(1 -2 2 -2 2 -1 1) $$
  $$ t^{-5} - 2t^{-4} + 2t^{-3} -2t^{-2} + 2t^{-1} - 1 + t $$
  $$ t=6, e^{x \log x} $$
  $$ \int y e^{-f}dV $$
  $$ \int(t^{-2} - t^{-1} + t - 1 + t)dt $$
  $$ = \int y e^{-f}dV $$
  $$ Li=6.941 Be=9.012 $$
  $$ Li - Be = 2.02988321$$
  $$ = \int y e^{-f}dV $$
  $$ {-2}(1 - 1 1 - 1 1) $$
  $$ \int e^{-x} \cdot x^{1 - t}dt $$
  $$ = \int (x - 1)^{t - 1}\cdot (t - 1)^{x - 1}dt $$
  $$ = e^{x \log x} $$
  $$ Li=6.941 - C=12.01 $$
  $$ 10, 3 2 1 1 1$$
  補空間の双曲体積
  $$ \text{第一周期 H    He} $$
  $$ \text{第二周期 Li Be    Ne} $$
  $$ {d \over df}F = 2{\int{{(R + \nabla_{i}\nabla_{j}f)} \over 
  {- (R + \Delta)}}}e^{-f}dV $$
  ８種類の組み合わせ

  これらより、質量差エネルギーをJones多項式で表せられることが、
  $$ {d \over df}F + \int C dx_m = \int \Gamma(\gamma)^{'}dx_m $$
  $$ = e^{-f} + e^{f} \ge e^{f} - e^{-f} $$
  $$ = 2(\cos(i x \log x) - i \sin(i x \log x) $$
  $$ = \beta(p,q) $$
  $$ = \int (R^{\mu\nu} + {1 \over 2}g_{ij}\Lambda) \mathrm{dvol} $$
  $$ = \int \kappa T^{\mu\nu} \mathrm{dvol} $$
  $$  = \int \Gamma(\gamma)^{'}dx_m $$
  とベータ関数から、一般相対性理論の多様体積分と渡せて、
  ガンマ関数における大域的積分多様体として、統一場理論が導けられる。

\end{document}
